\documentclass[12pt]{article}

\usepackage[margin=2cm]{geometry}
\usepackage{graphicx}
\usepackage{caption}
\usepackage[numbers]{natbib}
\usepackage{minted}   % syntax highlighting in code blocks
\usepackage{listings}
\usepackage{xcolor}
\usepackage[colorlinks=true, linkcolor=blue, urlcolor=blue, citecolor=blue]{hyperref}
\usepackage{microtype}  % nicer typesetting
\usepackage{libertinus}
\usepackage{inconsolata}  % nicer font for comments in code blocks


\linespread{1.05}

\title{\LARGE\textbf{Short Analysis Pipeline}\\\vspace{1mm}\large BIOLM0051}
\author{Archie Benn}
\date{}

% COMPILE WITH: pdflatex -shell-escape MAIN.tex 
\begin{document}
\maketitle

%%%%%%%%%%%%%%%%%%%%%%%%%%%%%%%%
\subsection*{Scenario}
A suitcase containing several packages of unidentified frozen meat has been intercepted at Heathrow Airport by UK Border Force officers during routine inspections. The shipment had no import documentation, raising concerns that the meat may originate from protected or endangered species.\newline

Given the serious implications for public health, food safety, and wildlife conservation, the case has been referred to the Animal and Plant Health Agency (APHA). As a bioinformatician assisting with the inquiry, you have been tasked with determining the species of origin for each meat sample. DNA barcoding has already been performed in the laboratory, and you have now received the resulting FASTAQ sequence files for analysis.

\section*{Abstract} %200 WORDS
\section*{Introduction}
\section*{Methods}

Some python code: 
\begin{minted}[fontsize=\small, bgcolor=gray!5, framesep=1.5mm, frame=single]{python}
from Bio import SeqIO
for record in SeqIO.parse("sequence.fasta", "fasta"):
    print(record.id, record.seq.translate())
\end{minted} 
\newline
Some Bash script:
\begin{minted}[fontsize=\small, frame=single, framesep=1.5mm, bgcolor=gray!5]{bash}
#!/bin/bash
#SBATCH --job-name=blast_job
#SBATCH --partition=teach_cpu
#SBATCH --time=01:00:00
#SBATCH --mem=2000M

blastn -db nt -query sequence.fa -out results.out -remote
\end{minted}
%%%%%%%%%%%%%%%%%%%%%%%%%%%%%%%%



%%%%%%%%%%%%%%%%%%%%%%%%%%%%%%%%
\section*{Results}
In this study, we analysed the gene expression levels across multiple samples using RNA-seq data. The expression values $x_i$ for gene $i$ were normalised using the log-transformation:
\[
y_i = \log_2(x_i + 1)
\]
to stabilise variance. Data processing was performed using Python scripts:

\begin{minted}[frame=single, framesep=1.5mm, bgcolor=gray!5]{python}
import pandas as pd
data = pd.read_csv("expression.csv")
normalized = data.apply(lambda x: np.log2(x+1))
\end{minted}
%%%%%%%%%%%%%%%%%%%%%%%%%%%%%%%%



%%%%%%%%%%%%%%%%%%%%%%%%%%%%%%%%
\section*{Discussion}
Yadadadada \citep{kent2002blat} and then also \citep{wolin1997microbe} and maybe even from this guy \cite{baker1997signaling,wolin1997microbe}
%%%%%%%%%%%%%%%%%%%%%%%%%%%%%%%%



%%%%%%%%%%%%%%%%%%%%%%%%%%%%%%%% check .bib ordering for years and surnames etc.
\bibliographystyle{unsrtnat}
\bibliography{references}
%%%%%%%%%%%%%%%%%%%%%%%%%%%%%%%%



%%%%%%%%%%%%%%%%%%%%%%%%%%%%%%%%
\section*{Supplementary Data}
Supplementary data can be found within the GitHub repository for this project: \url{https://github.com/archiebenn/intro_to_bioinformatics_report.git}
%%%%%%%%%%%%%%%%%%%%%%%%%%%%%%%%
\end{document}